% ==========================================================================
% Capítulo: Fases de desarrollo y cronograma
% Estado: se registra en ../STATUS.md, no aquí.
%
% Debe contener (project-context/formato-anteproyecto.md):
%   - FASES DE DESARROLLO DEL PROYECTO: en correspondencia con los objetivos
%     específicos (cadena de dependencias técnicas y duración estimada).
%   - CRONOGRAMA (diagrama de Gantt, ruta crítica e hitos).
% ==========================================================================
\chapter{Fases de desarrollo y cronograma}\label{cha:fases-cronograma}

Este capítulo define las fases de desarrollo del proyecto a partir de la cadena de dependencias técnicas que articula los objetivos específicos formulados en el \cref{cha:objetivos}, establece la duración estimada de cada una de ellas y presenta el cronograma general de ejecución mediante un diagrama de Gantt, señalando la distribución temporal de los \emph{sprints}, los hitos de entrega y la ruta crítica que rige la construcción de la plataforma.

\section{Fases de desarrollo del proyecto}
\label{sec:fases-desarrollo}

Las fases de desarrollo del proyecto guardan correspondencia directa con los objetivos específicos formulados en el \cref{cha:objetivos} y se estructuran según una cadena estricta de dependencias técnicas. El proceso inicia con la puesta en marcha del clúster sobre las estaciones de la sala 104M (\cref{sec:req-orquestacion}, grupo «Clúster base») y el despliegue de los motores de inferencia. Sobre esa infraestructura se instrumenta el módulo de observabilidad (\cref{sec:req-telemetria}), necesario para recolectar telemetría y métricas operativas antes de aplicar restricciones de recursos. Posteriormente, se implementa la gobernanza de cuotas dentro del clúster (\cref{sec:req-gobernanza}, grupos «Identidad y control de acceso», «Cuotas y aislamiento» y «Admisión, prioridad y reservas») y se construye la plataforma web de despliegue y administración, la cual expone la interfaz de despliegue del \cref{obj:orquestacion} y el módulo administrativo del \cref{obj:gobernanza} sobre las capacidades ya estabilizadas. El ciclo culmina con la evaluación empírica del sistema integrado mediante las pruebas de carga del grupo «Medición de desempeño» (\cref{sec:req-evaluacion}). La \cref{tab:fases-desarrollo} sintetiza esta secuencia con la duración estimada de cada fase en \emph{sprints} de tres semanas.

\begin{xltabular}{\textwidth}{p{0.2\textwidth} X p{0.14\textwidth} p{0.11\textwidth}}
        \caption{Fases de desarrollo del proyecto, su correspondencia con los objetivos específicos y su duración.}
        \label{tab:fases-desarrollo} \\
        \toprule
        \textbf{Fase} & \textbf{Actividades principales} & \textbf{Objetivo específico} & \textbf{Duración} \\
        \midrule
        \endfirsthead
        \toprule
        \textbf{Fase} & \textbf{Actividades principales} & \textbf{Objetivo específico} & \textbf{Duración} \\
        \midrule
        \endhead
        \bottomrule
        \endfoot
        Puesta en marcha del clúster &
        Grupos «Clúster base», «Orquestación de cargas» y «Motores y modelos» de la \cref{sec:req-orquestacion} &
        \Cref{obj:orquestacion} & 2 \emph{sprints} \\
        \addlinespace
        Observabilidad de infraestructura &
        Los cuatro grupos de la \cref{sec:req-telemetria} &
        \Cref{obj:telemetria} & 2 \emph{sprints} \\
        \addlinespace
        Gobernanza de cuotas &
        Grupos «Identidad y control de acceso», «Cuotas y aislamiento» y «Admisión, prioridad y reservas» de la \cref{sec:req-gobernanza} &
        \Cref{obj:gobernanza} & 2 \emph{sprints} \\
        \addlinespace
        Plataforma de despliegue y administración &
        Grupo «Interfaz de despliegue» de la \cref{sec:req-orquestacion} y grupo «Módulo administrativo» de la \cref{sec:req-gobernanza} &
        \Cref{obj:orquestacion}, \cref{obj:gobernanza} & 3 \emph{sprints} \\
        \addlinespace
        Evaluación empírica &
        Grupo «Medición de desempeño» de la \cref{sec:req-evaluacion} &
        \Cref{obj:evaluacion} & 1 \emph{sprint} \\
\end{xltabular}

\section{Cronograma}
\label{sec:cronograma}

El proyecto inicia el 10 de agosto de 2026 y finaliza el 31 de mayo de 2027. Entre el 10 de agosto y el 20 de septiembre de 2026 el equipo dedica una etapa previa a la formulación del anteproyecto, el levantamiento de requisitos con el laboratorio y la definición de la arquitectura propuesta. Los \emph{sprints} de desarrollo comienzan el 21 de septiembre de 2026 y se suceden en ciclos de tres semanas hasta el cierre del proyecto. Las actividades de cada \emph{sprint} se derivan de los grupos de requerimientos presentados en las \cref{sec:req-telemetria,sec:req-gobernanza,sec:req-orquestacion,sec:req-evaluacion}, y las asume el equipo del proyecto en su conjunto, razón por la cual el cronograma omite una columna de responsable. La \cref{fig:gantt} presenta el cronograma completo, con las fechas de cada \emph{sprint}, los grupos de requerimientos que atiende y el hito de entrega del anteproyecto. Como cada fase depende de que la anterior quede estable, las cinco fases forman una sola cadena que constituye en conjunto la ruta crítica del proyecto, de modo que el retraso de cualquier \emph{sprint} desplaza el cierre.

% [verify: ADR-033: el sprint 4 termina el 13-dic, y el sprint 5 (11-31 ene) inicia en enero]
La ejecución del proyecto se estructura en dos grandes etapas consecutivas. La primera etapa abarca de agosto a diciembre de 2026, periodo en el que se consolida la formulación del anteproyecto y se ejecutan los cuatro primeros \emph{sprints} de desarrollo, correspondientes a la puesta en marcha del clúster de cómputo y a la instrumentación de observabilidad de infraestructura. La segunda etapa comprende de enero a mayo de 2027, periodo durante el cual se implementan el sistema de gobernanza de cuotas, la plataforma de despliegue y administración, y la evaluación empírica del sistema, culminando con el informe final del proyecto.

\begin{landscape}
\begin{figure}[p]
  \centering
  \begin{tikzpicture}[font=\small, yscale=0.92]
    % --- Eje de meses ---------------------------------------------------
    \draw[black!25] (0.000,0.20) -- (0.000,-13.00);
    \node[anchor=south] at (0.869,0.20) {ago};
    \draw[black!25] (1.738,0.20) -- (1.738,-13.00);
    \node[anchor=south] at (2.579,0.20) {sep};
    \draw[black!25] (3.419,0.20) -- (3.419,-13.00);
    \node[anchor=south] at (4.288,0.20) {oct};
    \draw[black!25] (5.157,0.20) -- (5.157,-13.00);
    \node[anchor=south] at (5.998,0.20) {nov};
    \draw[black!25] (6.839,0.20) -- (6.839,-13.00);
    \node[anchor=south] at (7.708,0.20) {dic};
    \draw[black!25] (8.575,0.20) -- (8.575,-13.00);
    \node[anchor=south] at (9.444,0.20) {ene};
    \draw[black!25] (10.313,0.20) -- (10.313,-13.00);
    \node[anchor=south] at (11.099,0.20) {feb};
    \draw[black!25] (11.883,0.20) -- (11.883,-13.00);
    \node[anchor=south] at (12.752,0.20) {mar};
    \draw[black!25] (13.621,0.20) -- (13.621,-13.00);
    \node[anchor=south] at (14.461,0.20) {abr};
    \draw[black!25] (15.302,0.20) -- (15.302,-13.00);
    \node[anchor=south] at (16.171,0.20) {may};
    \draw[black!25] (17.040,0.20) -- (17.040,-13.00);
    \node[anchor=south] at (4.288,0.65) {\textbf{2026}};
    \node[anchor=south] at (12.808,0.65) {\textbf{2027}};
    \draw (0,0.00) -- (17.040,0.00);

    % --- Bandas de fase (tono por fase, ver leyenda), a la izquierda de
    % las etiquetas de fila, que a su vez quedan junto al eje -----------
    \fill[black!25] (-3.6,-1.0) rectangle (-3.1,-3.0);
    \node[rotate=90, font=\scriptsize] at (-3.35,-2.0) {F1};
    \fill[black!40] (-3.6,-3.0) rectangle (-3.1,-5.0);
    \node[rotate=90, font=\scriptsize] at (-3.35,-4.0) {F2};
    \fill[black!55] (-3.6,-6.0) rectangle (-3.1,-8.0);
    \node[rotate=90, font=\scriptsize] at (-3.35,-7.0) {F3};
    \fill[black!70] (-3.6,-8.0) rectangle (-3.1,-11.0);
    \node[rotate=90, font=\scriptsize] at (-3.35,-9.5) {F4};
    \fill[black!85] (-3.6,-11.0) rectangle (-3.1,-12.0);
    \node[rotate=90, font=\scriptsize, text=white] at (-3.35,-11.5) {F5};

    % --- Etiquetas de fila (número de \emph{sprint} y grupos que atiende) --
    \node[anchor=east, align=right, font=\scriptsize, text width=2.9cm] at (-0.2,-0.5) {Etapa previa};
    \node[anchor=east, align=right, font=\scriptsize, text width=2.9cm] at (-0.2,-1.5) {\emph{Sprint} 1\\Clúster base};
    \node[anchor=east, align=right, font=\scriptsize, text width=2.9cm] at (-0.2,-2.5) {\emph{Sprint} 2\\Orquestación y motores};
    \node[anchor=east, align=right, font=\scriptsize, text width=2.9cm] at (-0.2,-3.5) {\emph{Sprint} 3\\Telemetría (infra. y \emph{hardware})};
    \node[anchor=east, align=right, font=\scriptsize, text width=2.9cm] at (-0.2,-4.5) {\emph{Sprint} 4\\Telemetría de carga y paneles};
    \node[anchor=east, align=right, font=\scriptsize, text width=2.9cm] at (-0.2,-5.5) {Hito\\Entrega del anteproyecto};
    \node[anchor=east, align=right, font=\scriptsize, text width=2.9cm] at (-0.2,-6.5) {\emph{Sprint} 5\\Identidad y acceso};
    \node[anchor=east, align=right, font=\scriptsize, text width=2.9cm] at (-0.2,-7.5) {\emph{Sprint} 6\\Cuotas, admisión y reservas};
    \node[anchor=east, align=right, font=\scriptsize, text width=2.9cm] at (-0.2,-8.5) {\emph{Sprint} 7\\Interfaz de despliegue (servicio)};
    \node[anchor=east, align=right, font=\scriptsize, text width=2.9cm] at (-0.2,-9.5) {\emph{Sprint} 8\\Interfaz de despliegue (UI)};
    \node[anchor=east, align=right, font=\scriptsize, text width=2.9cm] at (-0.2,-10.5) {\emph{Sprint} 9\\Módulo administrativo};
    % [verify: ADR-031: diseño propio o reutilización del harness del laboratorio]
    \node[anchor=east, align=right, font=\scriptsize, text width=2.9cm] at (-0.2,-11.5) {\emph{Sprint} 10\\Medición de desempeño};
    \node[anchor=east, align=right, font=\scriptsize, text width=2.9cm] at (-0.2,-12.5) {Cierre};

    % --- Barras, con las fechas de inicio y fin junto a cada una -----------
    \fill[black!10] (0.504,-0.84) rectangle (2.858,-0.16);
    \node[anchor=south, font=\tiny] at (1.681,-0.12) {10 ago -- 20 sep};
    \fill[black!25] (2.858,-1.84) rectangle (4.036,-1.16);
    \node[anchor=south, font=\tiny] at (3.447,-1.12) {21 sep -- 11 oct};
    \fill[black!25] (4.036,-2.84) rectangle (5.213,-2.16);
    \node[anchor=south, font=\tiny] at (4.625,-2.12) {12 oct -- 1 nov};
    \fill[black!40] (5.213,-3.84) rectangle (6.390,-3.16);
    \node[anchor=south, font=\tiny] at (5.802,-3.12) {2 nov -- 22 nov};
    \fill[black!40] (6.390,-4.84) rectangle (7.567,-4.16);
    \node[anchor=south, font=\tiny] at (6.979,-4.12) {23 nov -- 13 dic};
    \node[draw, fill=black, diamond, inner sep=2pt] at (7.006,-5.5) {};
    \node[anchor=west, font=\tiny] at (7.15,-5.5) {Inicios de dic. 2026};
    \fill[black!55] (9.136,-6.84) rectangle (10.313,-6.16);
    \node[anchor=south, font=\tiny] at (9.725,-6.12) {11 ene -- 31 ene};
    \fill[black!55] (10.313,-7.84) rectangle (11.491,-7.16);
    \node[anchor=south, font=\tiny] at (10.902,-7.12) {1 feb -- 21 feb};
    \fill[black!70] (11.491,-8.84) rectangle (12.668,-8.16);
    \node[anchor=south, font=\tiny] at (12.080,-8.12) {22 feb -- 14 mar};
    \fill[black!70] (12.668,-9.84) rectangle (13.845,-9.16);
    \node[anchor=south, font=\tiny] at (13.257,-9.12) {15 mar -- 4 abr};
    \fill[black!70] (13.845,-10.84) rectangle (15.022,-10.16);
    \node[anchor=south, font=\tiny] at (14.434,-10.12) {5 abr -- 25 abr};
    \fill[black!85] (15.022,-11.84) rectangle (16.199,-11.16);
    \node[anchor=south, font=\tiny] at (15.611,-11.12) {26 abr -- 16 may};
    \fill[black!10] (16.199,-12.84) rectangle (17.040,-12.16);
    \node[anchor=south, font=\tiny] at (16.620,-12.12) {17 -- 31 may};

    % --- Leyenda de fases -----------------------------------------------
    \node[anchor=north west, font=\scriptsize, text width=17.5cm] at (-3.6,-13.35)
      {F1 Puesta en marcha del clúster \quad F2 Observabilidad de infraestructura \quad F3 Gobernanza de cuotas\\
       F4 Plataforma de despliegue y administración \quad F5 Evaluación empírica};
  \end{tikzpicture}
  \caption[Diagrama de Gantt del proyecto]{Diagrama de Gantt del proyecto. Cada barra indica sus fechas y, para los \emph{sprints}, el grupo de requerimientos que atiende, en el orden de las \cref{sec:req-telemetria,sec:req-gobernanza,sec:req-orquestacion,sec:req-evaluacion}; el tono lateral identifica la fase (leyenda F1--F5) y el rombo marca la entrega del anteproyecto a inicios de diciembre de 2026. Las cinco fases forman una sola cadena de dependencias y constituyen en conjunto la ruta crítica del proyecto. Elaboración propia.}
  \label{fig:gantt}
\end{figure}
\end{landscape}

